% !TeX root = ../main.tex
% !TeX program = pdflatex
% !TeX encoding = UTF-8
% Biên dịch tài liệu gốc bằng pdfLaTeX.

\chapter{Tối Ưu Hóa Đa Điểm Đến}
\label{chap:multilocation}

\section{Mô Tả Bài Toán Định Tuyến Đa Điểm}

Trong nhiều tình huống thực tế, một lộ trình phải đi qua nhiều địa điểm chứ không chỉ một điểm xuất phát và một điểm đích. Ví dụ: một chuyến giao hàng dừng tại nhiều địa chỉ hoặc chuỗi các công việc cần chạy trong một chuyến đi. Đây là bài toán kinh điển về việc tìm thứ tự ghé thăm sao cho tổng chi phí di chuyển là nhỏ nhất, còn gọi là bài toán Người Chào Hàng (Traveling Salesman Problem - TSP). Cho một tập hợp các địa điểm phải ghé thăm đúng một lần từ một điểm xuất phát cố định, nhiệm vụ là chọn thứ tự ghé thăm để giảm thiểu tổng chi phí của các đoạn đường đi qua.

Số lượng thứ tự ghé thăm có thể có tăng theo giai thừa của số lượng địa điểm, do đó việc kiểm tra mọi thứ tự là không khả thi nếu số điểm dừng quá lớn. Vì vậy, cần các phương pháp hợp lý tùy theo quy mô của bài toán.

\section{Giải Thích Các Phương Pháp Được Áp Dụng}

Trong đồ án này, nhóm đã triển khai hai phương pháp khác nhau để giải quyết bài toán đa điểm, tùy thuộc vào số lượng địa điểm cần đến:

\textbf{1. Duyệt Cạn (Brute Force TSP):} Được sử dụng khi số lượng điểm đến nhỏ. Thuật toán tạo ra tất cả các hoán vị của thứ tự ghé thăm. Chi phí đường đi ngắn nhất giữa từng cặp điểm được tính trước bằng Dijkstra. Sau đó, tổng chi phí của từng hoán vị được tính và thứ tự có chi phí thấp nhất được chọn. Phương pháp này đảm bảo tìm ra thứ tự tối ưu tuyệt đối, nhưng độ phức tạp tăng theo $O(N!)$; chương trình giới hạn tối đa tám đích.

\textbf{2. Láng Giềng Gần Nhất (Nearest Neighbor TSP):} Bắt đầu từ điểm xuất phát cố định, thuật toán liên tục chọn điểm chưa thăm có chi phí đường đi Dijkstra nhỏ nhất cho đến khi ghé đủ địa điểm. Phương pháp này nhanh hơn duyệt cạn nhưng không bảo đảm tối ưu toàn cục.

\section{So Sánh Kết Quả Tối Ưu Hóa}

Với bài toán xuất phát từ DL01 và đi qua các điểm DL03, DL07, DL15. Chi phí ban đầu nếu đi theo thứ tự nhập (DL01 $\rightarrow$ DL03 $\rightarrow$ DL07 $\rightarrow$ DL15) được so sánh với thứ tự được hệ thống tối ưu bằng thuật toán Brute Force.

\begin{table}[h!]
\centering
\caption{Thứ tự ban đầu so với thứ tự đã tối ưu}
\label{tab:multilocation_opt}
\begin{tabular}{lcc}
\toprule
Trường hợp & Thứ Tự Ghé Thăm & Tổng Chi Phí \\
\midrule
Thứ tự ban đầu & DL01 $\rightarrow$ DL03 $\rightarrow$ DL07 $\rightarrow$ DL15 & 2,785366 \\
Thứ tự đã tối ưu & DL01 $\rightarrow$ DL15 $\rightarrow$ DL03 $\rightarrow$ DL07 & 1,848521 \\
\bottomrule
\end{tabular}
\end{table}

Với kịch bản S0 và cấu hình \texttt{balanced}, duyệt cạn xét đủ sáu hoán vị
và giảm chi phí từ 2,785366 xuống 1,848521, tương đương \textbf{33,63\%}.

\section{Mức Độ Tối Ưu Của Kết Quả}

Tùy vào phương pháp được lựa chọn (dựa trên số lượng điểm dừng) mà kết quả trả về sẽ là tối ưu hay xấp xỉ:

\begin{itemize}[itemsep=2pt]
  \item \textbf{Duyệt Cạn (Brute Force TSP):} kết quả thu được là \textbf{tối ưu tuyệt đối} vì thuật toán đã vét cạn tất cả các hoán vị có thể có của lộ trình và đối chiếu chi phí của chúng.
  \item \textbf{Láng Giềng Gần Nhất (Nearest Neighbor TSP):} kết quả chỉ mang
  tính \textbf{xấp xỉ}. Với ví dụ trên, thuật toán chọn DL01 $\rightarrow$
  DL03 $\rightarrow$ DL15 $\rightarrow$ DL07 và đạt chi phí 1,920959, cao hơn
  nghiệm chính xác 1,848521 khoảng 3,92\%, nhưng không phải xét mọi hoán vị.
\end{itemize}
